\section{Proof Sketch}

The proof uses two coefficient pivots.  On $[-1,1]$, solve the root equation
for $\alpha_0$.  Lemma~\ref{lem:inner-derivative} bounds the derivative of
the resulting sweep, and Lemma~\ref{lem:inner-image-length} converts this
into an image-length estimate.  On each connected outer component, divide by
$\theta^{d-1}$ and solve for $\alpha_{d-1}$.
Propositions~\ref{prop:positive-outer-certificate} and
\ref{prop:negative-outer-certificate} establish the same certificate
separately on the two signs, without crossing zero or assigning $\pm1$ to an
outer chart.

The probabilistic interface requires some care because the hypotheses bound
only the means of random conditional-density caps.  Lemma~\ref{lem:kernel-caps}
constructs measurable, version-independent representatives of those caps and
proves simultaneous domination of all Borel pivot sets.  Lemma~\ref{lem:borel-sections}
proves that the root events are Borel and identifies their exact conditional
sections.  Proposition~\ref{prop:matching-disintegration} then matches each
section to the correct regular conditional kernel.

Propositions~\ref{prop:inner-chart-probability},
\ref{prop:positive-chart-probability}, and
\ref{prop:negative-chart-probability} integrate the deterministic image
bounds against their matching random caps.  Lemma~\ref{lem:three-piece-decomposition}
partitions the parameter interval exactly, and
Propositions~\ref{prop:weighted-chart-bound} and
\ref{prop:exact-chart-maximum} combine the three event controls into the
weighted estimate and then the exact maximum.  Finally,
Proposition~\ref{prop:class-supremum} takes the two suprema, while
Proposition~\ref{prop:fixed-eta-polynomial-specialization} gives the explicit
polynomial envelope.

The nonemptiness clause is independent of the root-hitting argument.  Three
degree-specific kernel propositions compute both witness caps as $1/(2R)$.
The remaining witness propositions verify cube support, ambient singularity,
and middle-coordinate dependence, and Lemma~\ref{lem:witness-threshold}
isolates the only use of the threshold $1/2$.  The final appendix proof
assembles these results into Theorem~\ref{thm:main}.
